\section{安全探索与长期运行}
\subsection{约束优化与软边界}
\begin{figure}[H]
\centering
\includegraphics[width=0.9\textwidth]{slides-images/slide-06.jpg}
\caption{Slide 6 将安全约束视为 reward penalty，用于 soft policy gradient。}
\end{figure}

\begin{knowledgebox}{安全探索策略}
\begin{itemize}
\item 将 safety constraint 转化为 penalty term $c(s_t,a_t)$；\\
\item 用 Lagrangian multiplier $\lambda$ 控制 penalty 强度；\\
\item 在 uncertainty 高时（如 new goal），降低 exploration rate，避免 high-risk 动作。
\end{itemize}
\end{knowledgebox}

\subsection{长时间 horizon 的监控}
Chelsea 强调：Autonomy 不仅是一次训练，而是多次 deployment 的 property，因此需要系统级监控：
\begin{itemize}
\item 记录每个 goal 的 success/failure 轨迹；
\item track $\Delta reward$ 与 $\Delta constraint$ 以判断 drifting；
\item 对于重复失败，动态回退 reward 模型或缩小 goal space。
\end{itemize}

\begin{warningbox}{部署时留有回退通道}
部署 pipeline 需保留“回滚 checkpoint”与“重新训练 reward 估计”的机制，以便在出现 hallucination reward、reset policy 失效等问题时快速救火。
\end{warningbox}

\subsection{本章小结}
安全探索通过 penalty/constraint 架构保障行为，长期运行则靠监控统计与回退计划，并最终形成可靠的 Autonomy 系统。
\newpage

